\documentclass[12pt]{article}
\usepackage[utf8]{inputenc}
\usepackage[T1]{fontenc}
\usepackage{hyperref}
\title{Middle English Language Resources for K-12 Curriculum Development}
\author{Compiled by Claude Sonnet 5.5 (claude-sonnet-5-5) for Andr\'as Kornai}
\date{}
\begin{document}
\maketitle

\section{Introduction}
Middle English (ISO 639-3 enm; Glottocode midd1317) is a historical stage of English, roughly 1175--1500, with no native speakers. A heritage curriculum would be a reading and philology programme built on surviving texts. The search found mainly university-level reference tools and nothing designed for K-12 classrooms.

\section{Primary Educational Resources}
No school primers or K-12 textbooks for Middle English turned up in the search. University course resource lists and a language-aids page may be adapted by teachers:
\begin{itemize}
\item University of Georgia, language aids page (comprehension, pronunciation, syntax): \url{https://ctcamp.franklinresearch.uga.edu/language-aids}
\item University of Kentucky course resource page: \url{https://www.uky.edu/~kiernan/ENG421/resource.htm}
\end{itemize}

\section{Linguistic and Reference Materials}
\begin{itemize}
\item Middle English Dictionary (MED), University of Michigan, listed by the Medieval Academy of America: \url{https://mdr-maa.org/resource/middle-english-dictionary/}
\item Middle English Compendium (MED with a bibliography and a corpus of Middle English prose and verse), library guide: \url{https://guides.lib.udel.edu/az/middle-english-compendium}
\item Lexilogos page linking the MED, the corpus and older glossaries: \url{https://www.lexilogos.com/anglais_moyen.htm}
\item Medieval-studies links list: \url{https://tigerweb.towson.edu/duncan/hellinks.html}
\item Mayhew, \emph{A Concise Dictionary of Middle English} (book listing only): \url{https://www.goodreads.com/book/show/64990807-a-concise-dictionary-of-middle-english}
\end{itemize}

\section{Religious and Community Resources}
None were found in the search.

\section{Digital and Audio Resources}
The online MED and Compendium above are the only digital resources found. No audio or video resources were identified in the search.

\section{Curriculum Development Context}
Little exists for school use. A curriculum would need to adapt university reference tools (dictionary, corpus) and create graded readers and pronunciation materials. No grammar textbook was identified in the search results.
\end{document}